% TikZ-схема фотостанда для гл. 7 § 7.3.
% Подключается через % TikZ-схема фотостанда для гл. 7 § 7.3.
% Подключается через % TikZ-схема фотостанда для гл. 7 § 7.3.
% Подключается через % TikZ-схема фотостанда для гл. 7 § 7.3.
% Подключается через \input{figures/diagrams/photo_stand.tex}.
\begin{tikzpicture}[
    every node/.style={font=\footnotesize},
    box/.style={draw, thick, fill=gray!8},
    phone/.style={draw, thick, fill=white, rounded corners=1.5pt,
                  minimum width=40mm, minimum height=10mm},
    epd/.style={draw, thick, fill=white, minimum width=34mm, minimum height=8mm},
    label/.style={font=\scriptsize, align=left, inner sep=1pt},
    sidelight/.style={draw=orange!80!black, thick, fill=yellow!50,
                      minimum width=4mm, minimum height=8mm},
]

% --- Корпус коробки (вид сбоку, разрез) -----------------------------
% Стенки и дно — верх открыт.
\draw[box]
    (-4.0, -2.8) -- (-4.0, 1.6)        % левая стенка
    (-4.0, -2.8) -- (4.0, -2.8)        % дно
    (4.0, -2.8) -- (4.0, 1.6);         % правая стенка

% Чёрное матовое дно (штриховка внутри)
\fill[pattern=north east lines, pattern color=black!70]
    (-3.85, -2.7) rectangle (3.85, -2.4);

% Жирные внутренние стенки = «чёрные стенки»
\draw[line width=1.8pt, black] (-3.85, -2.7) -- (-3.85, 1.45);
\draw[line width=1.8pt, black] ( 3.85, -2.7) -- ( 3.85, 1.45);

% --- Телефон лежит СВЕРХУ коробки, перекрывая отверстие -------------
% (объектив снизу телефона смотрит в открытое отверстие коробки)
\node[phone] (phone) at (0, 2.0) {};
\node[font=\scriptsize] at (0, 2.0) {Телефон (камерой вниз)};

% Объектив (снизу телефона, внутрь коробки)
\filldraw[fill=black!50] (0, 1.55) circle (1.8mm);
\node[label] at (0.55, 1.55) [right] {объектив};

% Отверстие в верхе коробки — рисуем разрыв стенок (уже сделано выше)
\node[label, align=center] at (-3.0, 2.4) {верх открыт\\(отверстие)};
\draw[->, thin, gray] (-2.55, 2.30) -- (-1.0, 1.65);

% --- Боковое освещение (телефон-фонарик с диффузором или лампа) -----
\node[sidelight, rotate=-30] (lampL) at (-3.0, 0.4) {};
\node[sidelight, rotate=30]  (lampR) at ( 3.0, 0.4) {};
\node[label, align=center] at (-3.0, 1.05) {боковой\\свет};
\node[label, align=center] at ( 3.0, 1.05) {боковой\\свет};

% Лучи света от боковых источников на EPD (диффузные пунктиры)
\draw[-{Stealth[length=1.5mm]}, dashed, orange!70!black, thick]
    (-2.7, 0.2) -- (-0.7, -2.05);
\draw[-{Stealth[length=1.5mm]}, dashed, orange!70!black, thick]
    ( 2.7, 0.2) -- ( 0.7, -2.05);

% --- EPD-панель в центре дна, лицом вверх ---------------------------
\node[epd] (epd) at (-0.6, -2.18) {EPD (Waveshare 2.13"V4)};

% --- Калибровочный квадрат рядом с EPD ------------------------------
\filldraw[draw=black, very thick, fill=white] (2.0, -2.35) rectangle (2.5, -1.85);
\draw[->, thin, gray] (2.65, -1.65) -- (2.40, -1.95);
\node[label, align=left] at (2.7, -1.45) [right] {калибр.\\квадрат\\(опц.)};

% --- Расстояние от объектива до EPD (двойная стрелка слева) ---------
\draw[<->, thick, red!80!black] (-3.4, 1.55) -- (-3.4, -2.18)
    node[midway, left, label, text=red!80!black] {15--20\,см};

% --- Подпись чёрного дна -------------------------------------------
\draw[->, thin, gray] (-3.2, -3.05) -- (-3.4, -2.55);
\node[label, align=left] at (-3.2, -3.18) {чёрное матовое дно};

% --- Заголовок вида ------------------------------------------------
\node[label] at (0, -3.55) [font=\scriptsize\itshape] {вид сбоку (разрез коробки)};

\end{tikzpicture}
.
\begin{tikzpicture}[
    every node/.style={font=\footnotesize},
    box/.style={draw, thick, fill=gray!8},
    phone/.style={draw, thick, fill=white, rounded corners=1.5pt,
                  minimum width=40mm, minimum height=10mm},
    epd/.style={draw, thick, fill=white, minimum width=34mm, minimum height=8mm},
    label/.style={font=\scriptsize, align=left, inner sep=1pt},
    sidelight/.style={draw=orange!80!black, thick, fill=yellow!50,
                      minimum width=4mm, minimum height=8mm},
]

% --- Корпус коробки (вид сбоку, разрез) -----------------------------
% Стенки и дно — верх открыт.
\draw[box]
    (-4.0, -2.8) -- (-4.0, 1.6)        % левая стенка
    (-4.0, -2.8) -- (4.0, -2.8)        % дно
    (4.0, -2.8) -- (4.0, 1.6);         % правая стенка

% Чёрное матовое дно (штриховка внутри)
\fill[pattern=north east lines, pattern color=black!70]
    (-3.85, -2.7) rectangle (3.85, -2.4);

% Жирные внутренние стенки = «чёрные стенки»
\draw[line width=1.8pt, black] (-3.85, -2.7) -- (-3.85, 1.45);
\draw[line width=1.8pt, black] ( 3.85, -2.7) -- ( 3.85, 1.45);

% --- Телефон лежит СВЕРХУ коробки, перекрывая отверстие -------------
% (объектив снизу телефона смотрит в открытое отверстие коробки)
\node[phone] (phone) at (0, 2.0) {};
\node[font=\scriptsize] at (0, 2.0) {Телефон (камерой вниз)};

% Объектив (снизу телефона, внутрь коробки)
\filldraw[fill=black!50] (0, 1.55) circle (1.8mm);
\node[label] at (0.55, 1.55) [right] {объектив};

% Отверстие в верхе коробки — рисуем разрыв стенок (уже сделано выше)
\node[label, align=center] at (-3.0, 2.4) {верх открыт\\(отверстие)};
\draw[->, thin, gray] (-2.55, 2.30) -- (-1.0, 1.65);

% --- Боковое освещение (телефон-фонарик с диффузором или лампа) -----
\node[sidelight, rotate=-30] (lampL) at (-3.0, 0.4) {};
\node[sidelight, rotate=30]  (lampR) at ( 3.0, 0.4) {};
\node[label, align=center] at (-3.0, 1.05) {боковой\\свет};
\node[label, align=center] at ( 3.0, 1.05) {боковой\\свет};

% Лучи света от боковых источников на EPD (диффузные пунктиры)
\draw[-{Stealth[length=1.5mm]}, dashed, orange!70!black, thick]
    (-2.7, 0.2) -- (-0.7, -2.05);
\draw[-{Stealth[length=1.5mm]}, dashed, orange!70!black, thick]
    ( 2.7, 0.2) -- ( 0.7, -2.05);

% --- EPD-панель в центре дна, лицом вверх ---------------------------
\node[epd] (epd) at (-0.6, -2.18) {EPD (Waveshare 2.13"V4)};

% --- Калибровочный квадрат рядом с EPD ------------------------------
\filldraw[draw=black, very thick, fill=white] (2.0, -2.35) rectangle (2.5, -1.85);
\draw[->, thin, gray] (2.65, -1.65) -- (2.40, -1.95);
\node[label, align=left] at (2.7, -1.45) [right] {калибр.\\квадрат\\(опц.)};

% --- Расстояние от объектива до EPD (двойная стрелка слева) ---------
\draw[<->, thick, red!80!black] (-3.4, 1.55) -- (-3.4, -2.18)
    node[midway, left, label, text=red!80!black] {15--20\,см};

% --- Подпись чёрного дна -------------------------------------------
\draw[->, thin, gray] (-3.2, -3.05) -- (-3.4, -2.55);
\node[label, align=left] at (-3.2, -3.18) {чёрное матовое дно};

% --- Заголовок вида ------------------------------------------------
\node[label] at (0, -3.55) [font=\scriptsize\itshape] {вид сбоку (разрез коробки)};

\end{tikzpicture}
.
\begin{tikzpicture}[
    every node/.style={font=\footnotesize},
    box/.style={draw, thick, fill=gray!8},
    phone/.style={draw, thick, fill=white, rounded corners=1.5pt,
                  minimum width=40mm, minimum height=10mm},
    epd/.style={draw, thick, fill=white, minimum width=34mm, minimum height=8mm},
    label/.style={font=\scriptsize, align=left, inner sep=1pt},
    sidelight/.style={draw=orange!80!black, thick, fill=yellow!50,
                      minimum width=4mm, minimum height=8mm},
]

% --- Корпус коробки (вид сбоку, разрез) -----------------------------
% Стенки и дно — верх открыт.
\draw[box]
    (-4.0, -2.8) -- (-4.0, 1.6)        % левая стенка
    (-4.0, -2.8) -- (4.0, -2.8)        % дно
    (4.0, -2.8) -- (4.0, 1.6);         % правая стенка

% Чёрное матовое дно (штриховка внутри)
\fill[pattern=north east lines, pattern color=black!70]
    (-3.85, -2.7) rectangle (3.85, -2.4);

% Жирные внутренние стенки = «чёрные стенки»
\draw[line width=1.8pt, black] (-3.85, -2.7) -- (-3.85, 1.45);
\draw[line width=1.8pt, black] ( 3.85, -2.7) -- ( 3.85, 1.45);

% --- Телефон лежит СВЕРХУ коробки, перекрывая отверстие -------------
% (объектив снизу телефона смотрит в открытое отверстие коробки)
\node[phone] (phone) at (0, 2.0) {};
\node[font=\scriptsize] at (0, 2.0) {Телефон (камерой вниз)};

% Объектив (снизу телефона, внутрь коробки)
\filldraw[fill=black!50] (0, 1.55) circle (1.8mm);
\node[label] at (0.55, 1.55) [right] {объектив};

% Отверстие в верхе коробки — рисуем разрыв стенок (уже сделано выше)
\node[label, align=center] at (-3.0, 2.4) {верх открыт\\(отверстие)};
\draw[->, thin, gray] (-2.55, 2.30) -- (-1.0, 1.65);

% --- Боковое освещение (телефон-фонарик с диффузором или лампа) -----
\node[sidelight, rotate=-30] (lampL) at (-3.0, 0.4) {};
\node[sidelight, rotate=30]  (lampR) at ( 3.0, 0.4) {};
\node[label, align=center] at (-3.0, 1.05) {боковой\\свет};
\node[label, align=center] at ( 3.0, 1.05) {боковой\\свет};

% Лучи света от боковых источников на EPD (диффузные пунктиры)
\draw[-{Stealth[length=1.5mm]}, dashed, orange!70!black, thick]
    (-2.7, 0.2) -- (-0.7, -2.05);
\draw[-{Stealth[length=1.5mm]}, dashed, orange!70!black, thick]
    ( 2.7, 0.2) -- ( 0.7, -2.05);

% --- EPD-панель в центре дна, лицом вверх ---------------------------
\node[epd] (epd) at (-0.6, -2.18) {EPD (Waveshare 2.13"V4)};

% --- Калибровочный квадрат рядом с EPD ------------------------------
\filldraw[draw=black, very thick, fill=white] (2.0, -2.35) rectangle (2.5, -1.85);
\draw[->, thin, gray] (2.65, -1.65) -- (2.40, -1.95);
\node[label, align=left] at (2.7, -1.45) [right] {калибр.\\квадрат\\(опц.)};

% --- Расстояние от объектива до EPD (двойная стрелка слева) ---------
\draw[<->, thick, red!80!black] (-3.4, 1.55) -- (-3.4, -2.18)
    node[midway, left, label, text=red!80!black] {15--20\,см};

% --- Подпись чёрного дна -------------------------------------------
\draw[->, thin, gray] (-3.2, -3.05) -- (-3.4, -2.55);
\node[label, align=left] at (-3.2, -3.18) {чёрное матовое дно};

% --- Заголовок вида ------------------------------------------------
\node[label] at (0, -3.55) [font=\scriptsize\itshape] {вид сбоку (разрез коробки)};

\end{tikzpicture}
.
\begin{tikzpicture}[
    every node/.style={font=\footnotesize},
    box/.style={draw, thick, fill=gray!8},
    phone/.style={draw, thick, fill=white, rounded corners=1.5pt,
                  minimum width=40mm, minimum height=10mm},
    epd/.style={draw, thick, fill=white, minimum width=34mm, minimum height=8mm},
    label/.style={font=\scriptsize, align=left, inner sep=1pt},
    sidelight/.style={draw=orange!80!black, thick, fill=yellow!50,
                      minimum width=4mm, minimum height=8mm},
]

% --- Корпус коробки (вид сбоку, разрез) -----------------------------
% Стенки и дно — верх открыт.
\draw[box]
    (-4.0, -2.8) -- (-4.0, 1.6)        % левая стенка
    (-4.0, -2.8) -- (4.0, -2.8)        % дно
    (4.0, -2.8) -- (4.0, 1.6);         % правая стенка

% Чёрное матовое дно (штриховка внутри)
\fill[pattern=north east lines, pattern color=black!70]
    (-3.85, -2.7) rectangle (3.85, -2.4);

% Жирные внутренние стенки = «чёрные стенки»
\draw[line width=1.8pt, black] (-3.85, -2.7) -- (-3.85, 1.45);
\draw[line width=1.8pt, black] ( 3.85, -2.7) -- ( 3.85, 1.45);

% --- Телефон лежит СВЕРХУ коробки, перекрывая отверстие -------------
% (объектив снизу телефона смотрит в открытое отверстие коробки)
\node[phone] (phone) at (0, 2.0) {};
\node[font=\scriptsize] at (0, 2.0) {Телефон (камерой вниз)};

% Объектив (снизу телефона, внутрь коробки)
\filldraw[fill=black!50] (0, 1.55) circle (1.8mm);
\node[label] at (0.55, 1.55) [right] {объектив};

% Отверстие в верхе коробки — рисуем разрыв стенок (уже сделано выше)
\node[label, align=center] at (-3.0, 2.4) {верх открыт\\(отверстие)};
\draw[->, thin, gray] (-2.55, 2.30) -- (-1.0, 1.65);

% --- Боковое освещение (телефон-фонарик с диффузором или лампа) -----
\node[sidelight, rotate=-30] (lampL) at (-3.0, 0.4) {};
\node[sidelight, rotate=30]  (lampR) at ( 3.0, 0.4) {};
\node[label, align=center] at (-3.0, 1.05) {боковой\\свет};
\node[label, align=center] at ( 3.0, 1.05) {боковой\\свет};

% Лучи света от боковых источников на EPD (диффузные пунктиры)
\draw[-{Stealth[length=1.5mm]}, dashed, orange!70!black, thick]
    (-2.7, 0.2) -- (-0.7, -2.05);
\draw[-{Stealth[length=1.5mm]}, dashed, orange!70!black, thick]
    ( 2.7, 0.2) -- ( 0.7, -2.05);

% --- EPD-панель в центре дна, лицом вверх ---------------------------
\node[epd] (epd) at (-0.6, -2.18) {EPD (Waveshare 2.13"V4)};

% --- Калибровочный квадрат рядом с EPD ------------------------------
\filldraw[draw=black, very thick, fill=white] (2.0, -2.35) rectangle (2.5, -1.85);
\draw[->, thin, gray] (2.65, -1.65) -- (2.40, -1.95);
\node[label, align=left] at (2.7, -1.45) [right] {калибр.\\квадрат\\(опц.)};

% --- Расстояние от объектива до EPD (двойная стрелка слева) ---------
\draw[<->, thick, red!80!black] (-3.4, 1.55) -- (-3.4, -2.18)
    node[midway, left, label, text=red!80!black] {15--20\,см};

% --- Подпись чёрного дна -------------------------------------------
\draw[->, thin, gray] (-3.2, -3.05) -- (-3.4, -2.55);
\node[label, align=left] at (-3.2, -3.18) {чёрное матовое дно};

% --- Заголовок вида ------------------------------------------------
\node[label] at (0, -3.55) [font=\scriptsize\itshape] {вид сбоку (разрез коробки)};

\end{tikzpicture}
